\newcommand{\ModuleID}{EL-II.3}
\newcommand{\ModuleTitle}{Infinitives and Indirect Statement}
\newcommand{\ModuleStage}{Stage II — Core Grammar}
\newcommand{\ModuleUnit}{II.3}
\newcommand{\ModuleOutcome}{Form the six ordinary infinitives and preserve subject, voice, and relative time in complementary and indirect-statement structures}
\newcommand{\ModulePrerequisites}{EL-II.2 mastery: control comparison, adverbs, and operational numerals}
\newcommand{\ModuleSessionOne}{Build the infinitive voice/time axis and contrast complementary and indirect-statement minimal pairs, then complete Worksheet II.9 as the guided application.}
\newcommand{\ModuleSessionTwo}{Analyze accusative subjects and predicate complements in context, convert direct to reported statements across time, and complete Worksheet II.12 independently.}
\newcommand{\ModulePrintedPractice}{Worksheets II.9 (guided Variant A) and II.12 (independent Variant D), with terminal keys}
\newcommand{\ModuleExtensionPractice}{Worksheets II.10 and II.11 remain optional remediation or delayed-recall variants}
\newcommand{\ModuleNext}{EL-II.4, Participles and the Ablative Absolute}
\newcommand{\ModuleMastery}{Form at least 90 percent of eighteen prompted infinitives, identify subject and relative time in ten unseen constructions with no more than one error, and convert three differently timed direct statements into accurate indirect statement.}
\newcommand{\ModuleDelayNote}{}
\newcommand{\ModuleReferenceFocus}{%
  \begin{itemize}
    \item present, perfect, and future active and passive infinitives;
    \item time simultaneous with, prior to, or subsequent to the governing verb;
    \item complementary, subject, object, and indirect-statement infinitive roles; and
    \item accusative subject and accusative predicate complement in indirect statement.
  \end{itemize}}
\newcommand{\ModuleContrast}{A complementary infinitive completes another verb and normally shares its understood subject; an indirect statement contains its own accusative subject. Infinitive tense is relative to the governing verb, not an absolute calendar label.}
\newcommand{\ModuleLessonSource}{\CourseRoot/shared/content/volume2/all}
\newcommand{\ModuleExerciseSource}{\CourseRoot/shared/exercises/volume2/all}
\newcommand{\ModuleWorkedBank}{II}
\newcommand{\ModuleExerciseBank}{II}
\newcommand{\ModuleWorksheetVariants}{A,D}
\newcommand{\ModuleScopeNote}{This packet covers the ordinary six-form system and indirect statement; exceptional infinitive constructions remain dictionary- and context-dependent.}
\newcommand{\ModuleEnrichment}{%
  \SubmoduleExtension
    {The infinitive system}
    {A two-dimensional voice/time axis prevents form labels from being reduced to one English infinitive.}
    {Place all six lesson infinitives for one regular verb on an active/passive by simultaneous/prior/subsequent grid. Under each, write its morphological components rather than a full translation.}
    {State the relative time represented by present, perfect, and future infinitives.}
    {Present infinitives are ordinarily simultaneous with the governing verb, perfect infinitives prior, and future infinitives subsequent; the active/passive axis separately identifies voice. Every cell must use the formation printed in the lesson.}
  \SubmoduleExtension
    {Reading the forms in context}
    {Complementary and indirect-statement minimal pairs can contain the same infinitive while assigning different subject structure.}
    {Take one infinitive from the lesson and diagram its complementary use beside an indirect statement. Box the governing verb, circle any accusative subject, and draw the understood-subject link where no new subject is expressed.}
    {Give the decisive structural difference between the two constructions.}
    {A complementary infinitive completes the governing verb and ordinarily shares its subject; indirect statement has an accusative subject belonging to the reported proposition. The presence and role of that accusative, not an English \English{to} or \English{that}, decides the analysis.}
  \SubmoduleExtension
    {Writing laboratory}
    {A time axis checks whether conversion to reported speech preserves chronology when the reporting verb changes.}
    {Use the lesson's three direct-time prompts. For each, mark event time relative to reporting time, then convert after a present reporting verb and after a past reporting verb without changing that relation.}
    {Explain why a perfect infinitive need not refer to an event remote from the speaker's present.}
    {Each conversion must preserve simultaneous, prior, or subsequent relation to its governing report. A perfect infinitive marks priority relative to that report; absolute remoteness is not encoded by the infinitive label alone.}
}
